\documentclass[12pt]{article}
\usepackage[a4paper,margin=1in]{geometry}
\usepackage[utf8]{inputenc}
\usepackage[T1]{fontenc}
\usepackage{amsmath,amssymb}
\usepackage{textcomp}
\usepackage{parskip}
\begin{document}
\begin{center}
{\Large \textbf{2009 STEP III}}
\end{center}

\subsection*{Page 1}
\begin{flushleft}
\bigskip\textbf{Section A: Pure Mathematics}\par
\bigskip\textbf{Question 1.} The points S, T , U and V have coordinates (s, ms), (t, mt), (u, nu) and (v, nv), respectively.\par
\[The lines SV and U T meet the line y = 0 at the points with coordinates (p, 0) and (q, 0),\]
respectively. Show that\par
(m - n)sv\par
\[p = ,\]
ms - nv\par
and write down a similar expression for q.\par
Given that S and T lie on the circle x2 + (y - c)2 = r2, find a quadratic equation satisfied by\par
s and by t, and hence determine st and s + t in terms of m, c and r.\par
\[Given that S, T , U and V lie on the above circle, show that p + q = 0.\]
\infty\par
\bigskip\textbf{Question 2.} (i) Let y = a xn , where the coefficients a are independent of x and are such that this\par
n n\par
\[n=0\]
X\par
series and all others in this question converge. Show that\par
\infty\par
y\par
\[= na\]
xn-1\par
,\par
n\par
\[n=1\]
X\par
and write down a similar expression for y .\par
\[Write out explicitly each of the three series as far as the term containing a^3.\]
\textbf{(ii)} It is given that y satisfies the differential equation\par
00 0 3\par
\[xy - y + 4x y = 0 .\]
By substituting the series of part (i) into the differential equation and comparing coef-\par
\[ficients, show that a^1 = 0.\]
\[Show that, for n > 4,\]
\[a n = - a n-4 ,\]
n(n - 2)\par
\[and that, if a^0 = 1 and a^2 = 0, then y = cos(x^2) .\]
\[Find the corresponding result when a^0 = 0 and a^2 = 1.\]
\end{flushleft}
\subsection*{Page 2}
\begin{flushleft}
\bigskip\textbf{Question 3.} The function f(t) is defined, for t 6= 0, by\par
t\par
\[f(t) = .\]
et - 1\par
\textbf{(i)} By expanding et, show that lim f(t) = 1 . Find f 0 (t) and evaluate lim f 0 (t) .\par
\[t\to0 t\to0\]
\textbf{(ii)} Show that f(t) + 1 t is an even function. [Note: A function g(t) is said to be even if\par
g(t) \equiv g(-t).]\par
\textbf{(iii)} Show with the aid of a sketch that et(1 - t) 6 1 and deduce that f 0 (t) 6= 0 for t 6= 0.\par
Sketch the graph of f(t).\par
\bigskip\textbf{Question 4.} For any given (suitable) function f, the Laplace transform of f is the function F defined by\par
\infty\par
\[F(s) = e\]
-stf(t)dt\par
\[(s > 0) .\]
\[\int\]
\textbf{(i)} Show that the Laplace transform of e\par
-btf(t),\par
\[where b > 0, is F(s + b).\]
\textbf{(ii)} Show that the Laplace transform of f(at), where a > 0, is a -1F( s ) .\par
a\par
\textbf{(iii)} Show that the Laplace transform of f (t) is sF(s) - f(0) .\par
\textbf{(iv)} In the case f(t) = sin t, show that F(s) = .\par
\[s^2 + 1\]
Using only these four results, find the Laplace transform of e\par
-pt\par
\[cos qt , where p > 0 and q > 0.\]
\end{flushleft}
\subsection*{Page 3}
\begin{flushleft}
\bigskip\textbf{Question 5.} The numbers x, y and z satisfy\par
\[x + y + z = 1\]
2 2 2\par
\[x + y + z = 2\]
3 3 3\par
\[x + y + z = 3 .\]
Show that\par
\[yz + zx + xy = - .\]
\[Show also that x2y + x2z + y2z + y2x + z2x + z2y = -1 , and hence that\]
\[xyz = .\]
Let S = xn + yn + zn . Use the above results to find numbers a, b and c such that the relation\par
n\par
\[S n+1 = aS n + bS n-1 + cS n-2 ,\]
holds for all n.\par
\bigskip\textbf{Question 6.} Show that eiβ - eiα = 2 sin 1 (β - α) for 0 < α < β < 2\pi . Hence show that\par
¯\par
¯ e\par
iα\par
- e\par
¯\par
¯\par
iβ\par
e\par
iγ\par
- e\par
iδ\par
+ e\par
iβ\par
- e\par
iγ\par
e\par
iα\par
- e\par
iδ\par
\[= e\]
iα\par
- e\par
iγ\par
e\par
iβ\par
- e\par
iδ\par
,\par
¯ ¯ ¯ ¯ ¯ ¯ ¯ ¯ ¯ ¯ ¯ ¯\par
\[where 0 < α <¯ β < γ <¯ ¯δ < 2\pi. ¯ ¯ ¯ ¯ ¯ ¯ ¯ ¯ ¯\]
Interpret this result as a theorem about cyclic quadrilaterals.\par
\end{flushleft}
\subsection*{Page 4}
\begin{flushleft}
\bigskip\textbf{Question 7.} (i) The functions f (x) are defined for n = 0, 1, 2, . . . , by\par
n\par
1 df (x)\par
n\par
\[f^0(x) =\]
\[1 + x^2\]
\[and f n+1(x) =\]
dx\par
.\par
\[Prove, for n > 1, that\]
(1 + x )f n+1(x) + 2(n + 1)xf\par
n\par
\textbf{(x)} + n(n + 1)f n-1(x) = 0 .\par
\textbf{(ii)} The functions P (x) are defined for n = 0, 1, 2, . . . , by\par
n\par
\[P (x) = (1 + x\]
2 )n+1\par
f (x) .\par
n n\par
\[Find expressions for P^0(x), P^1(x) and P^2(x).\]
\[Prove, for n > 0, that\]
dP (x)\par
2 n\par
P n+1(x) - (1 + x ) + 2(n + 1)xP\par
n\par
\textbf{(x)} = 0 ,\par
dx\par
and that P (x) is a polynomial of degree n.\par
n\par
\bigskip\textbf{Question 8.} Let m be a positive integer and let n be a non-negative integer.\par
\textbf{(i)} Use the result lim e\par
-mttn\par
\[= 0 to show that\]
t\to\infty\par
\[lim xm(ln x)n = 0 .\]
\[x\to0\]
By writing xx as ex ln x show that\par
\[lim xx = 1 .\]
\[x\to0\]
\textbf{(ii)} Let I = xm(ln x)ndx . Show that\par
n\par
\[\int\]
n + 1\par
\[I n+1 = - I n\]
m + 1\par
and hence evaluate I .\par
n\par
\textbf{(iii)} Show that\par
\[xxdx = 1 - 1 2 + 1 3 - 1 4 + \cdots .\]
2 3 4\par
\[\int\]
¡ ¢ ¡ ¢ ¡ ¢\par
\end{flushleft}
\subsection*{Page 5}
\begin{flushleft}
\bigskip\textbf{Section B: Mechanics}\par
\bigskip\textbf{Question 9.} A particle is projected under gravity from a point P and passes through a point Q. The angles\par
of the trajectory with the positive horizontal direction at P and at Q are θ and φ, respectively.\par
The angle of elevation of Q from P is α.\par
\textbf{(i)} Show that tan θ + tan φ = 2 tan α.\par
\textbf{(ii)} It is given that there is a second trajectory from P to Q with the same speed of projec-\par
tion. The angles of this trajectory with the positive horizontal direction at P and at Q\par
0 0\par
are θ and φ , respectively. By considering a quadratic equation satisfied by tan θ, show\par
0 0 0\par
\[that tan(θ + θ ) = - cot α. Show also that θ + θ = \pi + φ + φ .\]
\bigskip\textbf{Question 10.} A light spring is fixed at its lower end and its axis is vertical. When a certain particle P rests\par
on the top of the spring, the compression is d. When, instead, P is dropped onto the top\par
of the spring from a height h above it, the compression at time t after P hits the top of the\par
spring is x. Obtain a second-order differential equation relating x and t for 0 6 t 6 T , where\par
T is the time at which P first loses contact with the spring.\par
Find the solution of this equation in the form\par
\[x = A + B cos(ωt) + C sin(ωt) ,\]
where the constants A, B, C and ω are to be given in terms of d, g and h as appropriate.\par
Show that\par
\[T = d/g 2\pi - 2 arctan 2h/d .\]
p ³ p ´\par
\end{flushleft}
\subsection*{Page 6}
\begin{flushleft}
\bigskip\textbf{Question 11.} A comet in deep space picks up mass as it travels through a large stationary dust cloud. It is\par
subject to a gravitational force of magnitude Mf acting in the direction of its motion. When\par
it entered the cloud, the comet had mass M and speed V . After a time t, it has travelled a\par
distance x through the cloud, its mass is M (1 + bx), where b is a positive constant, and its\par
speed is v.\par
\textbf{(i)} In the case when f = 0, write down an equation relating V , x, v and b. Hence find an\par
expression for x in terms of b, V and t.\par
\textbf{(ii)} In the case when f is a non-zero constant, use Newton's second law in the form\par
\[force = rate of change of momentum\]
to show that\par
f t + V\par
\[v = .\]
1 + bx\par
Hence find an expression for x in terms of b, V , f and t.\par
Show that it is possible, if b, V and f are suitably chosen, for the comet to move with\par
constant speed. Show also that, if the comet does not move with constant speed, its\par
speed tends to a constant as t \to \infty.\par
\end{flushleft}
\subsection*{Page 7}
\begin{flushleft}
\bigskip\textbf{Section C: Probability and Statistics}\par
12 (i) Albert tosses a fair coin k times, where k is a given positive integer. The number of\par
\[heads he gets is X^1. He then tosses the coin X^1 times, getting X^2 heads. He then tosses\]
\[the coin X^2 times, getting X^3 heads. The random variables X^4, X^5, . . . are defined\]
\[similarly. Write down E(X^1).\]
\[By considering E(X^2 X^1 = x^1), or otherwise, show that E(X^2) =\]
k.\par
\infty\par
¯\par
Find E(X ). ¯\par
i\par
\[i=1\]
X\par
\textbf{(ii)} Bertha has k fair coins. She tosses the first coin until she gets a tail. The number of\par
\[heads she gets before the first tail is Y^1. She then tosses the second coin until she gets\]
\[a tail and the number of heads she gets with this coin before the first tail is Y^2. The\]
k\par
\[random variables Y^3, Y^4, . . . , Y\]
k\par
\[are defined similarly, and Y = Y\]
i\par
.\par
\[i=1\]
P\par
Obtain the probability generating function of Y , and use it to find E(Y ), Var (Y ) and\par
\[P(Y = r).\]
\bigskip\textbf{Question 13.} (i) The point P lies on the circumference of a circle of unit radius and centre O. The angle, θ,\par
between OP and the positive x-axis is a random variable, uniformly distributed on the\par
\[interval 0 6 θ < 2\pi. The cartesian coordinates of P with respect to O are (X, Y ). Find\]
the probability density function for X, and calculate Var (X).\par
Show that X and Y are uncorrelated and discuss briefly whether they are independent.\par
\textbf{(ii)} The points P (i = 1, 2, . . . , n) are chosen independently on the circumference of\par
i\par
the circle, as in part (i), and have cartesian coordinates (X , Y ). The point P has\par
i i\par
1 n 1 n\par
\[coordinates (X, Y ), where X = X and Y = Y . Show that X and Y are\]
i i\par
n n\par
\[i=1 i=1\]
uncorrelated. P P\par
Show that, for large n, P |X| 6 \approx 0.95 .\par
n\par
à r !\par
\end{flushleft}
\end{document}
